\documentclass[11pt]{article}
\usepackage[a4paper,margin=26mm]{geometry}
\usepackage{amsmath,amssymb,mathtools}
\usepackage{unicode-math}
\setmainfont{Latin Modern Roman}
\setmathfont{Latin Modern Math}
\usepackage{xurl,hyperref}
\hypersetup{hidelinks,pdftitle={UnsolvedMath Problem 10300066}}
\newcommand{\sref}[2]{\hyperlink{source:#1:#2}{[S#2]}}
\newcommand{\cataloguescope}[1]{#1}
\begin{document}
% UnsolvedMath ID 10300066; source catalogue code AMR-102-0066
\setcounter{section}{10300065}
\section{Miscellaneous — Question 15.1}\label{10300066}
{\small\sffamily UnsolvedMath ID 10300066 · Topology}
\subsection{Definitions and mathematical statement}

\paragraph{Shared notation.}$\mathbb N=\{1,2,\ldots\}$, and $\mathbb Z,\mathbb Q,\mathbb R,\mathbb C$ denote the integers, rationals, reals and complex numbers. Any other conventions are specified locally.


A surface lamination $\Lambda$ is a closed subset of $M$ locally homeomorphic to a product of a disk with a closed transverse set. It is essential when it has no sphere leaf or torus leaf bounding a solid torus, and its complementary regions are irreducible, have incompressible boundary, and have no compressing monogon. Such a monogon is a compressing disk with one boundary arc on a boundary leaf and the other along a transverse interval fiber. It is genuine if some complementary region is not an interval bundle. A sublamination is a closed union of leaves. A branched surface $B\subset M$ is a smoothly tangent $2$-complex with the standard local branching models. Its fibered neighborhood $N(B)$ is an interval bundle with the usual identifications over the branch locus. A lamination is carried by $B$ if it lies in $N(B)$ transverse to interval fibers, and fully carried if it meets every fiber. Splitting means cutting this fibered neighborhood along carried surface pieces and taking the induced branched surface. A coorientation orders the points of a transverse interval. A transverse set is well-ordered when every nonempty subset has a least point in that order. An ordinal is the order type of a well-ordered set; ordinary branched-surface weights use numerical weights satisfying matching rules along branches.
\paragraph{Statement recorded in the supplied source.}
Describe which cooriented surface laminations in $3$-manifolds can have well-ordered transverse spaces, and the possible order types and topological consequences. Develop a useful theory of branched surfaces with ordinal-valued weights, specifying appropriate matching and realization rules instead of assuming real addition applies unchanged to ordinals. \sref{10300066}{2}
\subsection{Short English statement}
Classify laminations with well-ordered transversals and develop a meaningful theory of ordinal weights on branched surfaces.
\subsection{Sources}
\begin{description}
\item[\hypertarget{source:10300066:1}{S1}] ULAM.AI and the UnsolvedMath Contributors, \emph{UnsolvedMath: A Curated Collection of Open Mathematics Problems}, record 10300066 (AMR-102-0066). CC BY 4.0. \url{https://huggingface.co/datasets/ulamai/UnsolvedMath}.
\item[\hypertarget{source:10300066:2}{S2}] Danny Calegari, \emph{Problems in foliations and laminations of 3-manifolds}, version 0.78 (2002), Question 15.1; Definitions 1.1--1.10 and local definitions and remarks. \url{https://arxiv.org/abs/math/0209081v1}.
\end{description}
\label{end:10300066}
\end{document}
